% problem-000318 1 EDTA配位与分子结构
\section*{1. EDTA配位与分子结构}
\textit{下列为分问。}

\subsection*{1-1}
EDTA是⼄⼆胺四⼄酸的英⽂名称的缩写，市售试剂是其⼆⽔合⼆钠盐。

\subsection*{1-1}
-1 画出EDTA⼆钠盐⽔溶液中浓度最⾼的阴离⼦的结构简式。

\subsection*{1-1}
-2 Ca $( \mathrm { E D T A } ) ^ { 2 - }$ −溶液可⽤于静脉点滴以排除体内的铅。写出这个排铅反应的化学⽅程式（⽤Pb2+ 表⽰铅）。1-1-3 能否⽤EDTA⼆钠盐溶液代替 $\mathrm { C a ( E D T A ) ^ { 2 } }$ −溶液排铅？为什么？

\subsection*{1-2}
氨和三氧化硫反应得到⼀种晶体，熔点 $2 0 5 ~ ^ { \circ } \mathrm { C }$ ，不含结晶⽔。晶体中的分⼦有⼀个三重旋转轴，有极性。画出这种分⼦的结构式，标出正负极。

\subsection*{1-3}
$\mathrm { N a _ { 2 } [ F e ( C N ) _ { 5 } ( N O ) }$ ]的磁矩为零，给出铁原⼦的氧化态。 $\mathrm { N a _ { 2 } [ F e ( C N ) _ { 5 } ( N O ) } \mathrm { \dot { \Omega } ] }$ ]是鉴定 $S ^ { 2 - }$ 的试剂，⼆者反应得到紫⾊溶液，写出鉴定反应的离⼦⽅程式。

\subsection*{1-4}
$\mathrm { C a S O _ { 4 } { \cdot } 2 H _ { 2 } O }$ 微溶于⽔，但在 $\mathrm { H N O _ { 3 } }$ (1 M)、 $\mathrm { H C l O _ { 4 } }$ (1 M)中可溶。写出能够解释 $\mathrm { C a S O _ { 4 } }$ 在酸中溶解的反应⽅程式。

\subsection*{1-5}
取质量相等的2份 $\mathrm { P b S O _ { 4 } }$ （难溶物）粉末，分别加⼊ $\mathrm { H N O _ { 3 } }$ (3 M)和 $\mathrm { H C l O _ { 4 } } \left( 3 \mathrm { M } \right)$ ，充分混合， $\mathrm { P b S O _ { 4 } }$ 在 $\mathrm { H N O _ { 3 } }$ 能全溶，⽽在 $\mathrm { H C l O _ { 4 } }$ 中不能全溶。简要解释PbSO_{4}在H $\mathrm { N O _ { 3 } }$ 中溶解的原因。

\subsection*{1-6}
X和Y在周期表中相邻。 $\mathrm { C a C O _ { 3 } }$ 与X的单质⾼温反应，⽣成化合物B和⼀种⽓态氧化物；B与Y的单质反应⽣成化合物C和X的单质；B⽔解⽣成D；C⽔解⽣成E，E⽔解⽣成尿素。确定B、C、D、E、X和$\mathbf { Y } _ { \circ }$
